

%%%%%%%%%%%%%%%
\begin{frame}[fragile]{La négation}

On veut les logements qui ne proposent pas de Ski. 

\vfill
La requête
suivante n'est pas la bonne solution.

\vfill

\begin{minted}[baselinestretch=.9,
               fontsize=\small,
               framesep=2mm]{sql}
select distinct nom
from  Logement, Activité
where code = codeLogement
and codeActivité != 'Ski'
\end{minted}

\vfill

Elle donne les logements où il existe une  activité autre que 'Ski'

\vfill

 On veut ici exprimer une condition (négative) 
sur un \alert{ensemble}. ``Il ne doit pas exister
d'activité Ski parmi celles du logement.''

\end{frame}

%%%%%%%%%%%%%%%
\begin{frame}[fragile]{La bonne formulation}

``Il ne doit pas exister
d'activité Ski parmi celles du logement.''
\vfill

C'est donc une condition existentielle (négative). 

\vfill

\begin{minted}[baselinestretch=.9,
               fontsize=\small,
               framesep=2mm]{sql}
select nom
from Logement 
where not exists (select ''
                 from Activité 
                 where code = codeLogement
                 and codeActivité = 'Ski')
\end{minted}

\vfill

Rappel\,: le \sqlkw{select} du bloc imbriqué ne sert à rien
\end{frame}


%%%%%%%%%%%%%%%
\begin{frame}[fragile]{Le \sqlkw{not in} marche aussi}

On peut toujours formuler le \sqlkw{(not) exists}
de manière très proche avec \sqlkw{(not) in}
\vfill


\begin{minted}[baselinestretch=.9,
               fontsize=\small,
               framesep=2mm]{sql}
select nom
from Logement
where code not in (select codeLogement
                   from Activité
                   where codeActivité = 'Ski')
\end{minted}

\vfill

Autre syntaxe possible\,: avec la différence ensembliste (``tous les logements moins
ceux où on fait du ski''). Peu pratique.
\end{frame}


%%%%%%%%%%%%%%%
\begin{frame}[fragile]{La quantification universelle}

Je sais exprimer une condition existentielle. Quid de la
quantification \alert{universelle}\,?

\vfill

``Je veux les voyageurs qui sont allés dans \alert{tous} les logements'' 

\vfill

\alert{SQL peut le faire}. Si je suis allé dans tous les logements,
\alert{il n'existe pas} de logement où \alert{je ne suis pas} allé.

\vfill

La quantification universelle s'exprime avec une double
quantification existentielle et la négation (Ouf...).
\end{frame}

\begin{frame}[fragile]{La solution (à étudier)}

Double imbrication, double quantification, double négation

\vfill

\begin{minted}[baselinestretch=.9,
               fontsize=\small,
               framesep=2mm]{sql}
select distinct v.prénom, v.nom
from Voyageur as v
where  not exists 
     (select '' from Logement as l
      where not exists 
           (select ''  from Séjour as s
            where l.code = s.codeLogement
            and   v.idVoyageur = s.idVoyageur))
\end{minted}

\vfill

Peu courant, mais montre la puissance du langage. 

\vfill

SQL couvre tout ce qui s'exprime dans la logique classique,
dite ''du premier ordre``. Cf. le complément sur le calcul.

\end{frame}

%%%%
\begin{frame}{Bilan}

Nous savons maintenant (presque) tout faire\,!
\vfill

\begin{redbullet}
 \item Exprimer toutes les requêtes dites ''positives'', sélection, projection
     et jointure.
  \item Exprimer la négation avec un \sqlkw{not exists}
  \item Et même exprimer la quantification universelle.
\end{redbullet}

\vfill

SQL, langage de logicien, pas d'informaticien

\vfill

Il reste à étudier quelques extensions pratiques.
\end{frame}

